\section{Что уже измерено и не работает}

Это тоже результат, и он объясняет, почему решение выглядит так, а не проще. Если скрининг настолько
информативен, зачем городить совместное правдоподобие? Нельзя ли подключить его в лоб? Проверены
три способа.

\begin{table}[htbp]\centering\small
\begin{tabularx}{\textwidth}{@{}Xrrrl@{}}
\toprule
Подход & ST-RAE & $\Delta$ к своей базе & $\rho$ & $\Delta\rho$ \\
\midrule
Скрининг как признак, в лоб & 0.921 & $+0.138$ & 0.566 & $+0.009$ \\
То же с честной вложенной схемой & 0.788 & $0.000$ & 0.559 & $+0.006$ \\
Два этапа: скрининг, затем пересчёт & 1.006 & $+0.212$ & 0.561 & $+0.010$ \\
Смесь двух маршрутов пополам & 0.803 & $+0.009$ & 0.587 & $+0.036$ \\
\bottomrule
\end{tabularx}
\caption{Разности считаются от базовой линии \emph{того же прогона}: 0.783, 0.788, 0.794, 0.794.
Расходятся они не случайно --- это три разных скрипта с тремя разными конфигурациями, см. ниже, ---
поэтому сравнивать строки между собой по абсолютной величине нельзя, только по разностям.}
\label{tab:neg}
\end{table}

\emph{Почему базовые линии не совпадают.} Три скрипта, давшие эту таблицу, строят разбиение по
фингерпринту в 1024 бита, тогда как абляционная сетка и все остальные числа документа --- по 2048.
При 1024 битах кластеризация даёт 4694 кластера против 4703, и вектор фолдов другой: дайджест
\texttt{d7b7ca15} против \texttt{2d93c198}. Поэтому ни одно число этой таблицы не сопоставимо
с 0.767 из раздела~\ref{sec:proto}: различаются не только модели, но и выборки, на которых
они отложены.

Между собой фолды у трёх скриптов совпадают, и расхождение базовых линий объясняется целиком
конфигурацией: шесть физико-химических признаков против восьми даёт 0.783 против 0.788, а
250 итераций бустинга против 200 --- 0.788 против 0.794. Оба различия детерминированы и
воспроизводятся при повторном запуске точно. Разброс в 0.011 --- это не шум прогона, а три
разные модели на общих фолдах.

\emph{Скрининг как признак, в лоб.} Обучаем вспомогательную модель предсказывать $\log_2$fc, её выход
подаём как признак в основную. Ошибка --- вспомогательный признак посчитан моделью, для которой
обучающие строки основной модели были внутренними. На них он точнее, чем окажется на отложенных,
основная модель учится на него опираться, и опора рассыпается там, где её проверяют: 0.783 против
0.921. На честной вложенной схеме выигрыш ровно нулевой. Разница между 0.783 и 0.921 целиком
артефакт процедуры.

\emph{Двухэтапная инверсия.} Предсказать $\log_2$fc, потом перевести в $\pic$ монотонной калибровкой.
Обязана выдать число там, где правдоподобие плоское: в зонах насыщения из раздела~\ref{sec:cens} калибровка
не определена, и любое значение оттуда одинаково согласуется с показанием. Результат: абсолютная шкала
разваливается, а ранговая информация в среднем сохраняется --- но только в среднем, см. ниже.

\emph{Смесь.} Полусумма двух маршрутов. Стоит $+0.009$ ST-RAE и покупает $+0.036$ Спирмена, причём ранг
растёт на всех четырёх ферментах. Это уже почти приемлемый обмен, но он не решает задачу, а балансирует
между двумя плохими решениями.

Общий вывод: информация в скрининге есть, но «ранг улучшается всегда» --- верно только для макро. Поферментно двухэтапная инверсия ранг роняет, причём на том ферменте, вокруг которого построен весь документ: на 2D6 Спирмен идёт с 0.398 до 0.362, на 3A4 с 0.749 до 0.747, и держится макро только за счёт 1A2 ($+0.048$) и 2C9 ($+0.028$). Смесь двух маршрутов --- единственный из трёх вариантов,
который поднимает ранг на всех четырёх.

Информация в скрининге, таким образом, есть, и вся проблема в том, как перевести её в абсолютную шкалу
$\pic$, не потеряв ранг там, где он нам нужнее всего. Этим занимается совместное правдоподобие через
общую латентную кривую.

Сюда же два отрицательных результата из предыдущих разделов: решающий слой под ST-RAE имеет потолок
0.015 макро и в реалистичной версии сейчас теряет 0.018; выигрыш механистического блока по ST-RAE
неотличим от нуля и поштучно, и на макро.
